\documentstyle[%


righttag%


,11pt%


,amssymb%


,russian%


,izvru%


]{amsart}





\newcommand{\tl}[3]{\medskip\noindent{\bf#1}\ #2 \dotfill #3 \par} 


\pagestyle{empty}


% \pagestyle{plain}


\begin{document}


% \setcounter{page}{80}


\par


\par


\par


\bigskip





\centerline{\Large СОДЕРЖАНИЕ}


\bigskip


\bigskip


%%%%%%%%%% Use \newline %%%%%%%%%%%











\tl{Асеев В.В., Кузин Д.Г., Тетенов А.В.}


{Углы между множествами и склейка\newline квазисимметрических


отображений в метрических пространствах}{3}


\tl{Гарифьянов Ф.Н., Туре Б.}


{Проблема обращения особого интеграла на круговом\newline секторе}{14}


\tl{Жуковский Е.С.}


{Нелинейное уравнение Вольтерра в банаховом


функциональном\newline пространстве}{17}


\tl{Иванов А.Г.}


{Динамическая система сдвигов и существование решения


задачи почти\newline периодической оптимизации}{29}


\tl{Исаенко Ю.Я.} 


{Структура периодических матриц-функций


второго порядка, имеющих\newline определитель, тождественно равный единице}{47}


\tl{Свиридюк Г.А., Тринеева И.К.}


{Сборка Уитни в фазовом пространстве уравнения\newline Хоффа}{54}


\tl{Степанянц С.А.}{Необходимые условия тауберова типа


для методов суммирования\newline Чезаро}{61}


\tl{Хайруллин Р.С.} 


{О задаче типа Геллерстедта для уравнения второго рода}{72}


\bigskip





\centerline{КРАТКИЕ~СООБЩЕНИЯ}


\medskip





\tl{Андреева Т.Н.} 


{Внутренняя геометрия сетей 


на гиперповерхности конформного\newline пространства}{78}








\vfill


\noindent\raisebox{2pt}{\copyright} Казанский госудаpственный унивеpситет


\end{document}


